\documentclass{beamer}

\usetheme{Padova}

\title{\textbf{Project 15}: Extend \textbf{Aexp} to include expressions whose evaluation may modify the current state}
\subtitle{Software Verification - Homework on Semantics}
\author{Gusatto Derek}
\date{2197482}



\begin{document}
	\maketitle
	\begin{frame}{Outline}
		\tableofcontents
	\end{frame}
    
%----------------------------------------------------
\section{Introduction}
%----------------------------------------------------

%%%%%%%%%%%%%%%%%%%%%%%%%%%%%%%%%%%%%%%%%%%%%%%%%%%%%%%%%%
\begin{frame}{Introduction}

In the seen \texttt{While}-semantics, arithmetic expressions are pure, i.e. the evaluation of variables produces a value without modifing any state:

\[
\mathcal{A}: \textbf{Aexp} \to \text{State} \to \mathbb{Z}\]
\[\mathcal{A}[\![n]\!]\sigma \Rightarrow  n
\]
\[\mathcal{A}[\![x]\!]\sigma \Rightarrow  \sigma(x) 
\]

However, real programming languages include operators such as:

\begin{center}
\texttt{++x \quad x++ \quad --x \quad x--}
\end{center}

These constructs introduce side effects during evaluation.
\note{
In the standard While-language, arithmetic expressions are evaluated using a semantic function
$\mathcal{A}: \textbf{Aexp} \to \text{State} \to \mathbb{Z}$.

This means that expressions like variables and constants simply read the state and produce a value. For example, evaluating a variable x just returns the value of x in the current state, without modifying memory state itself.

This assumption works well for theoretical languages, but it is not realistic when compared to actual programming languages such as C, Java, or C++, where expressions can have side effects.
    }
\end{frame}
%%%%%%%%%%%%%%%%%%%%%%%%%%%%%%%%%%%%%%%%%%%%%%%%%%%%%%%%%%
%%%%%%%%%%%%%%%%%%%%%%%%%%%%%%%%%%%%%%%%%%%%%%%%%%%%%%%%%%
\begin{frame}{Goal}

The goal is to extend semantics while preserving key properties:
\begin{itemize}
    \item Determinism
    \item Compositionality
    \item Operational clearity
\end{itemize}
\note{
The goal of this extension is not just to add new syntax, but to do so in a controlled semantic way.
In particular, we want to preserve three key properties:
\begin{itemize}

    \item Determinism, so that programs have a unique meaning
\item Compositionality, so that the meaning of complex expressions depends on their parts
\item Operational clarity, meaning that the order of evaluation and state changes is explicit
     
\end{itemize}

Achieving all three simultaneously is non-trivial once side effects are introduced.}
\end{frame}
%%%%%%%%%%%%%%%%%%%%%%%%%%%%%%%%%%%%%%%%%%%%%%%%%%%%%%%%%%

%----------------------------------------------------
\section{New Evaluation}
%----------------------------------------------------
%%%%%%%%%%%%%%%%%%%%%%%%%%%%%%%%%%%%%%%%%%%%%%%%%%%%%%%%%%
\begin{frame}{New Evaluation}

\begin{alertblock}{Key Observation}
Expression evaluation must return \alert{both} a value and an updated state.
\[\mathcal{A}: \textbf{Aexp} \to \text{State} \to (\mathbb{Z},\text{State})\]
\[\mathcal{A}[\![a]\!]\sigma \Rightarrow  (v, \sigma')
\]
\end{alertblock}

Meaning:
\begin{itemize}
    \item value $\mathcal{A}[\![a]\!]\sigma$ is produced;
    \item state may change to $\sigma'$.
    
\end{itemize}
Hence expression become state tranformers.
\note{
The key observation is that expression evaluation can no longer return just a value.
Since expressions may modify the state, their evaluation must return both a value and an updated state.

Formally, the semantic function for arithmetic expressions now maps an expression and a state to a pair consisting of an integer and a new state.
Intuitively, evaluating an expression produces a value, but it may also transform the state.

This means that arithmetic expressions become state transformers rather than pure computations.

}
\end{frame}
%%%%%%%%%%%%%%%%%%%%%%%%%%%%%%%%%%%%%%%%%%%%%%%%%%%%%%%%%%

%----------------------------------------------------
\section{Evaluation Order}
%----------------------------------------------------
%%%%%%%%%%%%%%%%%%%%%%%%%%%%%%%%%%%%%%%%%%%%%%%%%%%%%%%%%%
\begin{frame}{Evaluation order matters}

Without side-effects, easily:
\[a_1+a_2 = a_2 + a_1\]

But with side effects, for example with increments:

\[\texttt{x++} + \texttt{++x} \neq \texttt{++x}+\texttt{x++}\]

In fact the evaluation order affects both the final value and the final state.
\note{
Once expressions have side effects, evaluation order becomes crucial.
In the pure setting, expressions like $a_1 + a_2$ and $a_2 + a_1$ are equivalent.
But with side effects, this is no longer true.

For example, the expression $(x++) + (++x)$ does not have the same meaning as $(++x) + (x++)$.)

Both the final value and the final state may differ depending on the order of evaluation.
}
\end{frame}
%%%%%%%%%%%%%%%%%%%%%%%%%%%%%%%%%%%%%%%%%%%%%%%%%%%%%%%%%%
%%%%%%%%%%%%%%%%%%%%%%%%%%%%%%%%%%%%%%%%%%%%%%%%%%%%%%%%%%
\begin{frame}{Left-to-Right Evaluation}
We need to explicitly define:
\begin{alertblock}{Evaluation Strategy}
    Subexpressions are evaluated left-to-right.
\end{alertblock}


\begin{exampleblock}{Example}	
If a binary expression $a_1 \ + \ a_2$ is evaluated in state $\sigma$:
\begin{enumerate}
    \item Evaluate $a_1$ in $\sigma$ obtaining $(v_1,\sigma_1)$
\item Evaluate $a_2$ in $\sigma_1$ obtaining $(v_2,\sigma_2)$
\item Return $(v_1 \ + \ v_2, \sigma_2)$
\end{enumerate}

\end{exampleblock}
\note{
To avoid ambiguity, we must explicitly define an evaluation strategy.

In this project, we adopt a left-to-right evaluation order for subexpressions.
For a binary expression like $a_1 + a_2$, we first evaluate $a_1$ in the current state, obtaining a value and an intermediate state.

Then we evaluate $a_2$ in that intermediate state.
Finally, we combine the two values and return the final state produced by evaluating $a_2$.
}
\end{frame}
%%%%%%%%%%%%%%%%%%%%%%%%%%%%%%%%%%%%%%%%%%%%%%%%%%%%%%%%%%
%%%%%%%%%%%%%%%%%%%%%%%%%%%%%%%%%%%%%%%%%%%%%%%%%%%%%%%%%%
\begin{frame}{Left-to-Right Evaluation}


Reasons:
\begin{itemize}
    \item Avoid semantic ambiguity;
    \item Guaranteee determinism;
    \item Match real languages (e.g. Java).
\end{itemize}
Without fixed order in evaluation:
\begin{itemize}
    \item multiple possible results;
    \item non deterministic semantics.
\end{itemize}
\note{
Fixing a left-to-right evaluation order has several advantages.

First, it avoids semantic ambiguity and ensures determinism.
Without a fixed order, the same program could have multiple meanings.

Second, it matches the behavior of many real programming languages, such as Java.

Overall, this choice makes the semantics both precise and realistic.
}
\end{frame}
%%%%%%%%%%%%%%%%%%%%%%%%%%%%%%%%%%%%%%%%%%%%%%%%%%%%%%%%%%

%----------------------------------------------------
\section{Semantics Domains}
%----------------------------------------------------
%%%%%%%%%%%%%%%%%%%%%%%%%%%%%%%%%%%%%%%%%%%%%%%%%%%%%%%%%%
\begin{frame}{Semantics Domains}

We follow the standard \texttt{While}-language setting.

\begin{align*}
Var & : \text{variables}\\
State &= Var \rightarrow \mathbb{Z}\\
\sigma & : \text{program state}
\end{align*}

Functional update:

\[
\sigma[x \mapsto v]
\]

represents a new state where variable $x$ has value $v$.

\note{
We keep the standard semantic domains of the While-language.
A state is a function from variables to integers.
We also use functional update notation:
$\sigma[x \mapsto v]$ represents a new state where variable x is updated to value v, and all other variables remain unchanged.

}
\end{frame}
%%%%%%%%%%%%%%%%%%%%%%%%%%%%%%%%%%%%%%%%%%%%%%%%%%%%%%%%%%

%----------------------------------------------------
\section{Extended Syntax}
%----------------------------------------------------
%%%%%%%%%%%%%%%%%%%%%%%%%%%%%%%%%%%%%%%%%%%%%%%%%%%%%%%%%%
\begin{frame}{Extended Syntax}

Arithmetic expressions:

\[
a ::= n \mid x \mid a_1 + a_2 \mid ++x \mid x++ \mid --x \mid x--
\]
\vspace{.5em}

\begin{block}{Note:}
			\begin{itemize}
			    \item Increment applies only to variables;
                \item Boolean expressions and Statements syntax remain unchanged, semantics is modified.
			\end{itemize}
		\end{block}

\note{
We extend the syntax of arithmetic expressions to include increment and decrement operators.

These operators apply only to variables, not to arbitrary expressions.

Importantly, the syntax of boolean expressions and statements remains unchanged.

However, even though their syntax is the same, their semantics will change because they now rely on arithmetic expressions with side effects.
}
\end{frame}
%%%%%%%%%%%%%%%%%%%%%%%%%%%%%%%%%%%%%%%%%%%%%%%%%%%%%%%%%%

%----------------------------------------------------
\section{\textbf{Aexp} Semantics}
%----------------------------------------------------
%%%%%%%%%%%%%%%%%%%%%%%%%%%%%%%%%%%%%%%%%%%%%%%%%%%%%%%%%%
\begin{frame}{Basic arithmetic semantics}

\[\mathcal{A}: \textbf{Aexp} \to \text{State} \to (\mathbb{Z},\text{State})\]
\vspace{.5em}
\[\mathcal{A}[\![n]\!]\sigma \Rightarrow (\mathcal{Z}[\![n]\!], \sigma) \]
\[\mathcal{A}[\![x]\!]\sigma \Rightarrow (\sigma(x), \sigma) \]
\[\frac{\mathcal{A}[\![a_1]\!] \sigma \Rightarrow (v_1, \sigma_1) \quad \mathcal{A}[\![a_2]\!] \sigma_1 \Rightarrow (v_2, \sigma_2)}{\mathcal{A}[\![a_1 \texttt{ bap } a_2]\!]\sigma \Rightarrow  (v_1 \texttt{ bap } v_2, \sigma_2)}\]
\[\texttt{bap} \in \{+, -, *\}\]

\note{
Let’s now look at the updated semantics for basic arithmetic expressions.

Constants and variables behave as before, except that they now explicitly return the unchanged state.

For binary operators, evaluation follows the left-to-right strategy we discussed earlier.

The key point is that the state produced by evaluating the left operand is passed to the right operand.}
\end{frame}
%%%%%%%%%%%%%%%%%%%%%%%%%%%%%%%%%%%%%%%%%%%%%%%%%%%%%%%%%%
%%%%%%%%%%%%%%%%%%%%%%%%%%%%%%%%%%%%%%%%%%%%%%%%%%%%%%%%%%
\begin{frame}{New Expressions}

\[\mathcal{A}: \textbf{Aexp} \to \text{State} \to (\mathbb{Z},\text{State})\]
\vspace{.5em}
\[\mathcal{A}[\![++x]\!]\sigma \Rightarrow (\sigma(x)+1, \sigma[x \mapsto \sigma(x)+1]) \]
\[\mathcal{A}[\![x++]\!]\sigma \Rightarrow (\sigma(x), \sigma[x \mapsto \sigma(x)+1]) \]
\vspace{.5em}
\[\mathcal{A}[\![--x]\!]\sigma \Rightarrow (\sigma(x)-1, \sigma[x \mapsto \sigma(x)-1]) \]
\[\mathcal{A}[\![x--]\!]\sigma \Rightarrow (\sigma(x), \sigma[x \mapsto \sigma(x)-1]) \]
\note{
This slide shows the semantics of the new expressions.
For pre-increment, we first increment the variable in the state and return the updated value.

For post-increment, we return the old value but still update the state afterward.

The same pattern applies symmetrically to pre-decrement and post-decrement.

These rules precisely capture the intuitive behavior of these operators in real languages.
}
\end{frame}
%%%%%%%%%%%%%%%%%%%%%%%%%%%%%%%%%%%%%%%%%%%%%%%%%%%%%%%%%%

%----------------------------------------------------
\section{\textbf{Bexp} Semantics}
%----------------------------------------------------
%%%%%%%%%%%%%%%%%%%%%%%%%%%%%%%%%%%%%%%%%%%%%%%%%%%%%%%%%%
\begin{frame}{Changes in the \textbf{Bexp} Semantics}

\[
\mathcal{B}: \textbf{Bexp} \to \text{State} \to (\mathbb{T}, \text{State})\]

\[\mathcal{B}[\![true]\!]\sigma \Rightarrow (tt, \sigma) \quad \mathcal{B}[\![false]\!]\sigma \Rightarrow (ff, \sigma)\]
\[\frac{\mathcal{A}[\![a_1 ]\!]\sigma \Rightarrow (v_1, \sigma_1) \quad \mathcal{A}[\![a_2 ]\!]\sigma_1 \Rightarrow (v_2, \sigma_2)}{\mathcal{B}[\![a_1 = a_2]\!]\sigma\Rightarrow(v_1 = v_2, \sigma_2)}\]

\[\frac{\mathcal{A}[\![a_1 ]\!]\sigma \Rightarrow (v_1, \sigma_1) \quad \mathcal{A}[\![a_2 ]\!]\sigma_1 \Rightarrow (v_2, \sigma_2)}{\mathcal{B}[\![a_1 \leq a_2]\!]\sigma\Rightarrow(v_1 \leq v_2, \sigma_2)}\]

\note{
Boolean expressions must also be updated, because they may contain arithmetic expressions with side effects.

As a result, boolean evaluation now returns both a boolean value and an updated state.

Equality and comparison operators evaluate their operands left-to-right, threading the state through the evaluation.

This ensures consistency with the arithmetic semantics.
}

\end{frame}
%%%%%%%%%%%%%%%%%%%%%%%%%%%%%%%%%%%%%%%%%%%%%%%%%%%%%%%%%%
%%%%%%%%%%%%%%%%%%%%%%%%%%%%%%%%%%%%%%%%%%%%%%%%%%%%%%%%%%
\begin{frame}{Short-Circuit Boolean Semantics}
%Since expressions may have side effects, evaluation order for logical operators must be explicitly defined.
%We adopt \textbf{short-circuit evaluation} (as in Java/C-like languages).
\[
\frac{
\mathcal B[\![b_1]\!]\sigma \Rightarrow (ff,\sigma_1)
}{
\mathcal B[\![b_1 \land b_2]\!]\sigma \Rightarrow (ff,\sigma_1)
}
\]
\[
\frac{
\mathcal B[\![b_1]\!]\sigma \Rightarrow (tt,\sigma_1)
\quad
\mathcal B[\![b_2]\!]\sigma_1 \Rightarrow (bv,\sigma_2)
}{
\mathcal B[\![b_1 \land b_2]\!]\sigma \Rightarrow (bv,\sigma_2)
}
\]

\[
\frac{
\mathcal B[\![b_1]\!]\sigma \Rightarrow (tt,\sigma_1)
}{
\mathcal B[\![b_1 \lor b_2]\!]\sigma \Rightarrow (tt,\sigma_1)
}
\]
\[
\frac{
\mathcal B[\![b_1]\!]\sigma \Rightarrow (ff,\sigma_1)
\quad
\mathcal B[\![b_2]\!]\sigma_1 \Rightarrow (bv,\sigma_2)
}{
\mathcal B[\![b_1 \lor b_2]\!]\sigma \Rightarrow (bv,\sigma_2)
}
\]



\begin{alertblock}{Short-circuiting}
Prevents unnecessary evaluation and avoids unintended side effects.
\end{alertblock}
\note{
Short-circuit semantics become especially important in the presence of side effects.

For conjunction, if the first operand evaluates to false, the second operand is not evaluated.

Similarly, for disjunction, if the first operand evaluates to true, the second operand is skipped.

This behavior prevents unnecessary evaluation and avoids unintended side effects, which closely matches real programming languages.
}
\end{frame}
%%%%%%%%%%%%%%%%%%%%%%%%%%%%%%%%%%%%%%%%%%%%%%%%%%%%%%%%%%

%%%%%%%%%%%%%%%%%%%%%%%%%%%%%%%%%%%%%%%%%%%%%%%%%%%%%%%%%%

%----------------------------------------------------
\section{\textbf{Stm} Semantics}
%----------------------------------------------------
%%%%%%%%%%%%%%%%%%%%%%%%%%%%%%%%%%%%%%%%%%%%%%%%%%%%%%%%%%
\begin{frame}{Changes in \textbf{Stm} Semantics}

The general structure of the statement semantics remains unchanged:

\[
\mathcal{S} : \textbf{Stm} \rightarrow (State \hookrightarrow State)
\]

However, since arithmetic and boolean expressions may now
produce side effects, all rules involving expressions must
thread the updated state produced during evaluation
\note{
The overall structure of statement semantics remains the same.

Statements still map states to states.

However, any rule that involves expression evaluation must now take into account the updated state returned by that evaluation.

In other words, states must be carefully threaded through all semantic rules.
}
\end{frame}
%%%%%%%%%%%%%%%%%%%%%%%%%%%%%%%%%%%%%%%%%%%%%%%%%%%%%%%%%%

%----------------------------------------------------
\section{Structural Operation Semantics}
%----------------------------------------------------
%%%%%%%%%%%%%%%%%%%%%%%%%%%%%%%%%%%%%%%%%%%%%%%%%%%%%%%%%%
\begin{frame}{Assignment}
Assignment must use the updated state returned by expression evaluation.

\[
\frac{\mathcal{A}[\![a]\!]\sigma \Rightarrow (v,\sigma_1)}{\langle x:=a,\sigma\rangle \Rightarrow \sigma_1[x \mapsto v]}
\]

Meaning:
\begin{itemize}
\item evaluate expression $a$ producing value $v$ and state $\sigma_1$;
\item update variable $x$ in the resulting state.
\end{itemize}
\note{
Assignment is a good example of this change.

First, the expression on the right-hand side is evaluated, producing a value and an intermediate state.

Then the variable on the left-hand side is updated in that intermediate state.

This ensures that side effects inside the expression are not lost.
}
\end{frame}

%%%%%%%%%%%%%%%%%%%%%%%%%%%%%%%%%%%%%%%%%%%%%%%%%%%%%%%%%%
%%%%%%%%%%%%%%%%%%%%%%%%%%%%%%%%%%%%%%%%%%%%%%%%%%%%%%%%%%
\begin{frame}{Conditional and Loop Semantics}

Boolean evaluation may modify the state.

\textbf{If statement:}

\[
\frac{\mathcal{B}[\![b]\!]\sigma \Rightarrow (tt,\sigma_1)}
{\langle \text{if } b \text{ then } S_1 \text{ else } S_2,\sigma \rangle
\Rightarrow
\langle S_1,\sigma_1\rangle}
\]\[
\frac{\mathcal{B}[\![b]\!]\sigma \Rightarrow (ff,\sigma_1)}
{\langle \text{if } b \text{ then } S_1 \text{ else } S_2,\sigma \rangle
\Rightarrow
\langle S_2,\sigma_1\rangle}
\]

\textbf{While statement:}

\[
\langle \text{while }b \text{ do }S, \sigma \rangle \Rightarrow \langle \text{if } b \text{ then } (S; \text{ while } b \text{ do } S) \text{ else } skip, \sigma \rangle
\]

\note{
Conditionals and loops are also affected, because boolean expressions may now modify the state.

In an if statement, the branch is chosen based on the boolean value, but the state produced by evaluating the condition is passed to the selected branch.

The while loop is defined in terms of the if statement, so it naturally inherits this behavior.
}
\end{frame}

%%%%%%%%%%%%%%%%%%%%%%%%%%%%%%%%%%%%%%%%%%%%%%%%%%%%%%%%%%

%----------------------------------------------------
\section{Conclusion}
%----------------------------------------------------
%%%%%%%%%%%%%%%%%%%%%%%%%%%%%%%%%%%%%%%%%%%%%%%%%%%%%%%%%%
\begin{frame}{Semantic Consequences}
\begin{exampleblock}{}
\textbf{Initial state:} $\qquad \qquad \qquad \sigma(x)=5$

\begin{columns}

\column{0.40\textwidth}
\textbf{Case 1:} $(--x)+(x++)$

\begin{itemize}
\item Evaluate $--x$:
\[
(4,\sigma_1), \quad \sigma_1(x)=4
\]
\item Evaluate $x++$ in $\sigma_1$:
\[
(4,\sigma_2), \quad \sigma_2(x)=5
\]
\end{itemize}
\[
\textbf{Result: } 4+4=8
\]

\column{0.40\textwidth}
\textbf{Case 2:} $(x++)+(--x)$

\begin{itemize}
\item Evaluate $x++$:
\[
(5,\sigma_1), \quad \sigma_1(x)=6
\]
\item Evaluate $--x$ in $\sigma_1$:
\[
(5,\sigma_2), \quad \sigma_2(x)=5
\]
\end{itemize}
\[
\textbf{Result: } 5+5=10
\]

\end{columns}
\end{exampleblock}


%\vspace{0.5em}

\begin{alertblock}{Evaluation order is semantically relevant}
Same final state, different result
\end{alertblock}
\note{
This example illustrates why evaluation order matters.

Starting from the same initial state, two expressions that differ only in evaluation order produce different results.

Interestingly, they may even end up in the same final state while producing different values.

This clearly shows that side effects make expression evaluation semantically richer and more subtle.
}
\end{frame}
%%%%%%%%%%%%%%%%%%%%%%%%%%%%%%%%%%%%%%%%%%%%%%%%%%%%%%%%%%
%%%%%%%%%%%%%%%%%%%%%%%%%%%%%%%%%%%%%%%%%%%%%%%%%%%%%%%%%%
\begin{frame}{Summary of the Extension}

We extended the While-language by introducing expressions
with side effects:

\begin{center}
\texttt{++x \quad x++ \quad --x \quad x--}
\end{center}

Key semantic changes:

\begin{itemize}
\item Arithmetic and boolean expressions are no longer pure;
\item Expression evaluation returns both a value and an updated state:
\[\mathcal{A}: \textbf{Aexp} \to \text{State} \to (\mathbb{Z},\text{State})\]
\item Evaluation order must be explicitly fixed
(left-to-right) to preserve determinism;
\item Statement semantics must propagate the state
produced during expression evaluation.
\end{itemize}
\note{
To conclude, we extended the While-language to support arithmetic expressions with side effects.

This required changing the semantics of arithmetic and boolean expressions so that they return both a value and a state.

We fixed a left-to-right evaluation order to preserve determinism.

Finally, we adapted statement semantics to correctly propagate state changes.

Overall, this extension brings the While-language closer to real programming languages while remaining formal and precise.
}
\end{frame}
%%%%%%%%%%%%%%%%%%%%%%%%%%%%%%%%%%%%%%%%%%%%%%%%%%%%%%%%%%


\end{document}
